\documentclass[10pt,a4paper]{amsart}
\usepackage[margin=19mm]{geometry}
\usepackage{amsmath,amssymb,mathtools}
\usepackage[scr=rsfs,cal=euler]{mathalfa}
\usepackage{xfrac}
\usepackage[unicode,hidelinks,bookmarks=false]{hyperref}
\newcommand{\ie}{i.e.\ }
\newcommand{\cf}{cf.\ }
\newcommand{\resp}{respectively\ }
\newcommand{\Subsection}{\subsection}
\providecommand{\nobreakdash}{\nobreak\hbox{-}}
\setlength{\emergencystretch}{2em}
\multlinegap0pt
\usepackage{amssymb}
\usepackage{tikz}
\newtheorem{cor}{Corollary}
\newtheorem{prop}{Proposition}
\newtheorem{thm}{Theorem}
\newcommand {\MWITT}   {\ensuremath{\operatorname{MWITT}}}
\newcommand {\Sign}    {{\ensuremath{\operatorname{Sign}}}}
\newcommand {\Witt}   {{\ensuremath{\operatorname{Witt}}}}
\newcommand {\cplx}  {\ensuremath{\mathbb{C}}}
\newcommand {\intg}  {\ensuremath{\mathbb{Z}}}
\newcommand {\pt}    {\ensuremath{\operatorname{pt}}}
\newcommand {\rat}   {\ensuremath{\mathbb{Q}}}
\newcommand {\tr}    {{\ensuremath{\operatorname{tr}}}}
\title{Equivariant L-Classes of Atiyah-Singer-Zagier Type for Singular Spaces}
\author{Markus Banagl}
\date{Source: arXiv:2604.14913v1 (CC BY 4.0)}
\begin{document}
\maketitle

\noindent\emph{Source and licence.} Equivariant L-Classes of Atiyah-Singer-Zagier Type for Singular Spaces. Version arXiv:2604.14913v1 (CC BY 4.0), distributed by arXiv under the Creative Commons Attribution 4.0 International licence, http://creativecommons.org/licenses/by/4.0/, as recorded in the arXiv licence metadata for this exact version and shown on the abstract page https://arxiv.org/abs/2604.14913v1. Full licence notice: \texttt{licenses/arxiv-2604.14913-CC-BY-4.0.txt}.\par\smallskip

\noindent\textit{Editorial scope.} Counted unit: the article's thm thm.freeaction (source0.tex lines 3726--3730). The article's complete original proof of it is delivered with it (source0.tex lines 3731--3867, one \texttt{proof} environment of 137 lines). No mathematics is written, repaired or re-derived: the statement and the proof are the article's own lines, in the article's own notation, and no retained line is substituted or re-flowed.  Retained uncounted as the source-local context the proof needs: lines 2936--2940; lines 3429--3432; lines 3504--3515.  Nothing was omitted from any retained span, so the package carries no omission record.  No retained line contains a citation, so the wrapper carries no bibliography and no bibliography entry.  No retained line contains a figure environment, an image-inclusion command or drawing code, so no image is delivered and the package declares no asset dependency.  Editorial formatting only, in addition: an AMS \texttt{amsart} preamble that restates the article's own declarations and macros used by the retained lines, namely the environments \texttt{cor}, \texttt{prop}, \texttt{thm} on separate counters, together with the standard packages those lines need.  The retained environments print in this document as Corollary 1, Proposition 1, Theorem 1 in order of appearance, whereas the article prints the same items with the numbering of its own full document.  The complete native source of the article as distributed in the arXiv e-print for this exact version is bundled unchanged as \texttt{sources/198565/source0.tex}.
\begin{cor} \label{cor.orbitspaceiswitt}
Let $G$ be a finite group and $X$ an oriented triangulated Witt pseudomanifold
upon which $G$ acts by orientation preserving simplicial maps.
Then the orbit space $X/G$ is an oriented Witt pseudomanifold.
\end{cor}

\begin{prop} \label{prop.wittgbordisminvariance}
(Banagl-Leichtnam-Piazza.) The equivariant signature $\Sign (g,X)$
of closed oriented Witt pseudomanifolds $X$ is a simplicial Witt $G$-bordism invariant.
\end{prop}

\begin{prop} \label{prop.evaloflgxissigngpreim}
Let $u\in H^i (S^i;\intg)$ be the generator compatible with the orientation
and let $\langle -,- \rangle$ denote evaluation of a cohomology class on
a homology class.
The equivariant class $L_* (g,X)$ satisfies
\begin{equation} \label{equ.fulgxissigngfinvp}
\langle f^* (u),~ L_* (g,X) \rangle = \Sign (g,f^{-1} (p)). 
\end{equation}
For $\dim X \leq 2(i-1),$
it is the unique $G$-invariant homology class satisfying this equation
for all $G$-invariant homotopy classes $f: X \to S^i$.
\end{prop}

\begin{thm} \label{thm.freeaction}
(Free action.)
Suppose that $G$ acts freely on $X$.
If $g\not= 1$, then $L_* (g,X) =0.$
\end{thm}

\begin{proof}
By Proposition \ref{prop.evaloflgxissigngpreim}, it suffices to show
that $\Sign (g,X)=0$ for free $2m$-dimensional
$G$-Witt spaces $X$, $g\in G - \{ 1 \}$.
Since $\Sign (g,X)$ depends only on the action of the subgroup 
$\langle g \rangle \cong \intg/_k$ generated by $g\not= 1,$
we may as well assume $G = \langle g \rangle$.
Let $BG$ be the corresponding classifying space.

Suppose first that $X$ has the form $X= Y \times \intg/_k$,
where $G=\intg/_k$ acts trivially on the Witt space $Y$ and by
translation on the second factor.
Then $IH_m (X) = \bigoplus_{j=0}^{k-1} IH_m (Y \times g^j)$
and the intersection form $B$ on $IH_m (X)$ is given by an orthogonal sum
of $k$ copies of the intersection form $B_Y$ on $IH_m (Y)$.
Choose a positive definite inner product $\langle \cdot, \cdot \rangle$
on $IH_m (Y) = IH_m (Y\times g^0)$ and extend it to $IH_m (X)$
by 
$\langle (v, g^j), (w, g^j) \rangle = \langle (v,g^0), (w,g^0)  \rangle$
and $\langle (v, g^i), (w, g^j) \rangle = 0$ for $g^i \not= g^j.$
This is then a positive definite and $G$-invariant inner product
on $IH_m (X)$. Suppose that $m$ is even.
Let $IH_+ (Y)$ and $IH_- (Y)$ be the positive and negative eigenspaces
of the operator $A_Y$ defined on $IH_m (Y)$ by 
$B_Y (v,w) = \langle v, A_Y w \rangle$, yielding a decomposition
$IH_m (Y) = IH_+ (Y) \oplus IH_- (Y)$.
The operator $A$ on $IH_m (X)$
defined by $B (v,w) = \langle v, A w \rangle$
is a block diagonal sum of $k$ copies of $A_Y$. 
The positive and negative eigenspaces of $A = k\cdot A_Y$
are given by
$IH_\pm = \bigoplus_j IH_\pm (Y) \times g^j$.
Now let $\{ (e^+_1, g^0),\ldots,  (e^+_s, g^0) \}$ be a basis for $IH_+ (Y)$.
Then a basis for $IH_+$ is given by
$\{ (e^+_1, g^j),\ldots,  (e^+_s, g^j) ~|~ j=0,\ldots, k-1  \}.$
If $g_*|_{IH_+}$ is written as a matrix with respect to this basis,
then this matrix has zero blocks, corresponding to $g^j$, along the
diagonal since $g_*$ acts by sending a basis vector $(e^+_i, g^j)$ to 
the basis vector $(e^+_i, g^{j+1})$.
This implies that the trace $\tr (g_*|_{IH_+})$ is zero.
Similarly,  $\tr (g_*|_{IH_-})=0$. Thus the difference $\Sign (g,X)$
of these two traces is zero.
The case of odd $m$ is treated similarly:
The operators $A_Y$ and $A$ are skew-adjoint and
the complex structure $J = A (AA^*)^{-1/2}$ on $IH_m (X)$
is the block diagonal sum of $k$ copies of the complex structure
$J_Y = A_Y (A_Y A_Y^*)^{-1/2}$ on $IH_m (Y)$.
Thus there is a direct sum decomposition of complex vector spaces
$(IH_m (X), J) = \bigoplus_{j=0}^{k-1} (IH_m (Y \times g^j), J_Y).$
Let $\{ (e_1, g^0),\ldots,  (e_s, g^0) \}$ be a basis for the complex
vector space $(IH_+ (Y), J_Y)$.
Then a basis for the complex vector space $(IH_m (X), J)$ is given by
$\{ (e_1, g^j),\ldots,  (e_s, g^j) ~|~ j=0,\ldots, k-1  \}.$
If the complex linear automorphism
$g_*|_{(IH_m (X), J)}$ is written as a matrix with respect to this basis,
then this matrix again has zero blocks along the
diagonal since $g_*|_{(IH_m (X), J)}$ acts by sending a basis vector $(e_i, g^j)$ to 
the basis vector $(e_i, g^{j+1})$.
This implies that the trace $\tr (g_*|_{(IH_m (X), J)})$ is zero.
Thus $\Sign (g,X)=0$ also when $m$ is odd.

We show next that some multiple of a free simplicial $G=\intg/_k$-Witt pseudomanifold $X$
is simplicially Witt $G$-bordant to a $G$-Witt pseudomanifold of the above form 
$Y \times \intg/_k$. This will imply the desired statement by the
Witt $G$-bordism invariance of the $g$-signature,
Proposition \ref{prop.wittgbordisminvariance}.
If $G=\langle g \rangle$ acts freely on $X$, then the orbit map
$\pi: X \to X/G$ is a regular covering space classified by a map
$f: Y := X/G \to BG$. The orbit space $X/G$ is an oriented Witt pseudomanifold by
Corollary \ref{cor.orbitspaceiswitt}.
Hence, the map $f$ defines an element
$[f] \in \Omega^\Witt_* (B\intg/_k)$, where $\Omega^\Witt_* (-)$ denotes
Witt bordism theory.
We recall that a connected space $S$ is called (homotopy theoretically) 
\emph{simple} if $\pi_1 (S)$ is abelian and the action of
$\pi_1 (S)$ on $\pi_n (S)$ is trivial for every $n\geq 2$.
A simple space has a rational localization $S_\rat$, which is again a simple space.
Standard localization methods in stable homotopy theory show that
$\widetilde{E}_n (S_\rat) = \widetilde{E}_n (S) \otimes \rat$
for every spectrum $E$.
The classifying space $BG=B\intg/_k = K(\intg/_k, 1)$ is a simple
space, since $\intg/_k$ is abelian and $\pi_n (K(\intg/_k, 1))=0$ for all
$n\geq 2$. (Explicitly, $B\intg/_k$ is an infinite dimensional lens space
$L = L^\infty_k$.)
The localization $L_\rat$ is homotopy equivalent to a point, since
$\pi_1 (L) \otimes \rat = \intg/_k \otimes \rat =0$,
and for $n\geq 2$, $\pi_n (L) \otimes \rat =0$.
Given any spectrum $E$, we deduce that
$\widetilde{E}_n (L) \otimes \rat = \widetilde{E}_n (L_\rat) 
= \widetilde{E}_n (\pt) \otimes \rat =0.$
The long exact sequence of the pair $(L, \{ l_0 \}),$ where $l_0 \in L$ is a base point,
shows that the inclusion $j:\{ l_0 \} \hookrightarrow L$
induces an isomorphism
$E_n (\{ l_0 \}) \otimes \rat \cong
   E_n (L) \otimes \rat.$
Applying this to the spectrum $E=\MWITT$ representing Witt bordism theory,
we receive an isomorphism
$\Omega^\Witt_n (\pt) \otimes \rat \cong
   \Omega^\Witt_n (B\intg/_k) \otimes \rat.$
We return to the element
$[f] \in \Omega^\Witt_* (BG)$ represented by the classifying map $f:X/G \to BG$.
The above isomorphism shows that some multiple of $[f]$ is
Witt bordant over $BG$ to a map which factors as
\[ Y \stackrel{c}{\longrightarrow} 
    \pt \stackrel{j}{\hookrightarrow}
     BG, \]
where $Y$ is some closed oriented Witt pseudomanifold.     
Let $F:W \to BG$ be such a bordism.     
Let $EG \to BG$ be the universal principal $G$-bundle.
(Explicitly, $EG = S^\infty$.) This is the universal cover of $BG=K(\intg/_k, 1)$.
Since $f$ classifies the $\intg/_k$-space $X$, we have
$X = f^* EG$. 
Let $V := F^* EG$ and $X' := c^* j^* EG$.
Then $V$ is a covering space of $W$ and thus also a Witt pseudomanifold.
The group $G=\intg/_k$ is the deck transformation group and in particular
operates freely on $V$.
Let $C$ be any simplicial complex triangulating $W$ which is fine enough 
so that every simplex
is contained in a connected open set which is evenly covered by $V\to W$.
Then $C$ can be lifted to a triangulation of $V$ by a simplicial complex $K$
such that $V\to W$ becomes a simplicial map $K\to C$
under the triangulation homeomorphisms.
Since the restriction of this map to any slice over an evenly covered subset
is invertible, the group of
deck transformations (i.e. $G$) acts simplicially on $K$.
Since $W$ is oriented, its covering space $V$ is oriented as well.
The boundary of $V$ is 
$(F|_{\partial W})^* EG = f^* EG \sqcup c^* j^* EG = X \sqcup X'.$
Thus $V$ is a simplicial Witt $G$-bordism between a multiple $rX$ of $X$
and $X'$. Now $j^* EG = \intg/_k,$ the principal $G$-bundle over a point.
Thus
$X' = c^* (\intg/_k) = Y \times \intg/_k,$
where $\intg/_k$ acts trivially on $Y$ and by translation on the second factor.
Hence by Proposition \ref{prop.wittgbordisminvariance}, 
$r\Sign (g,X) = \Sign (g, rX) = \Sign (g,X') = \Sign (g,Y\times \intg/_k)=0 \in \cplx$.
This implies that $\Sign (g,X)=0$.
\end{proof}

\end{document}
